\textbf{RBX1 is hard.} It is 108 amino acids long, half of which are
intrinsically disordered, and the folded RING-H2 core depends on three
structural Zn²⁺ ions held in a cross-brace coordination
{[}@zheng2002scf{]}. In vivo it is the catalytic tip of every
cullin-RING ligase (CRL), recruiting an E2 ubiquitin-conjugating enzyme
via its α2 helix and handing ubiquitin onto a substrate organised by the
rest of the complex {[}@petroski2005crl{]}. There is no published
designed binder against RBX1, no widely-used reference protein binder
beyond CUL1 itself, and no obvious druggable pocket --- the RING surface
is shallow, the tail is disordered, and the catalytic helix is small.

That is why we picked it. The GEM × Adaptyv 2026 competition follows two
prior community campaigns at Adaptyv (EGFR 2024 {[}@adaptyv2024egfr{]},
Nipah-G 2025--26 {[}@adaptyv2026nipah{]}) by escalating the target
difficulty. Submitters could enter any modality --- short miniproteins,
peptides, single-domain antibodies, scFvs --- provided they kept ≤250
aa, ≤75\% identity to known proteins, and ≤100 sequences per team. The
submission window opened on 15 Feb 2026 and closed on 26 Mar 2026.
Validation ran at the Adaptyv Foundry through April; results were
announced at the ICLR GEM workshop on 26 Apr 2026.

We report the wet-lab outcome of the 322-design batch here. The batch
intentionally mixes methods (RFdiffusion, BoltzGen, Mosaic, BindCraft 2,
MoPPIt, LFM2, hallucination on AF2, and 25 more) and modalities so that
no single design philosophy is privileged. The full submission table,
every ESMFold prediction, and every replicate SPR sensorgram is
published at ProteinBase under the ODC-ODbL licence.
